پروژه \lr{Now in Android} به‌صورت ماژولار سازمان‌دهی شده و ماژول‌ها معمولاً در سه دسته‌ی اصلی قرار می‌گیرند: \lr{app}، \lr{core} و \lr{feature}. ماژول \lr{app} نقش نقطه‌ی ورود را دارد؛ یعنی راه‌اندازی اپ، پیکربندی \lr{DI}، مسیریابی/ناوبری، اتصال صفحه‌ها به یکدیگر و تعریف وابستگی‌ها در سطح برنامه عمدتاً در این ماژول انجام می‌شود و خودِ منطق قابلیت‌ها تا حد امکان در آن قرار نمی‌گیرد. ماژول‌های \lr{core} بخش‌های زیرساختی و مشترک را فراهم می‌کنند (مانند \lr{core:network}، \lr{core:database}، \lr{core:datastore}، \lr{core:model} و \lr{core:designsystem}) و به‌عنوان سرویس‌های عمومی در اختیار سایر بخش‌ها قرار می‌گیرند تا قابلیت استفاده‌مجدد و کپسوله‌سازی حفظ شود. در مقابل، ماژول‌های \lr{feature} هر کدام یک قابلیت/صفحه‌ی مستقل از اپ (مثل علاقه‌مندی‌ها، بوکمارک‌ها، جستجو و …) را شامل می‌شوند و معمولاً \lr{UI} و منطق ارائه‌ی همان دامنه را نگه می‌دارند و برای دریافت داده به \lr{core} (یا \lr{core:data}) وابسته‌اند؛ این تفکیک باعث می‌شود توسعه، تست، و نگه‌داری هر قابلیت مستقل‌تر و مقیاس‌پذیرتر انجام شود.


\subsectionaddtolist{۱. چرا تمام ماژول‌های کاربردی اپ در ماژول \lr{app} ادغام نشده‌اند؟}

ادغام کردن تمام کدها در یک ماژول \lr{app} در پروژه‌های کوچک ممکن است قابل قبول باشد، اما در پروژه‌های واقعی و بزرگ باعث افزایش پیچیدگی و هزینه نگه‌داری می‌شود. دلایل اصلی ماژولار بودن در پروژه \lr{NiA} عبارت‌اند از:

\begin{itemize}
  \item \textbf{افزایش سرعت ساخت و کامپایل:}
  با جدا کردن قابلیت‌ها در ماژول‌های مستقل، Gradle می‌تواند فقط ماژول‌های تغییر یافته را بازسازی کند و حتی بخشی از کارها را موازی‌سازی کند. این موضوع در پروژه‌های بزرگ بسیار تاثیرگذار است.

  \item \textbf{\lr{Separation of Concerns}:}
  هر ماژول مسئول یک دامنه مشخص است مثلاً \lr{feature:search} یا
  \linebreak
  \lr{feature:bookmarks} و این باعث می‌شود کدها کمتر به هم وابسته شوند و از ایجاد \lr{Spaghetti Code} جلوگیری شود.

  \item \textbf{کاهش وابستگی‌ها و کپسوله‌سازی:}
  ماژول‌ها می‌توانند API مشخصی ارائه دهند و جزئیات داخلی خود را مخفی نگه دارند. در نتیجه تغییرات داخلی یک ماژول، کمترین اثر را روی سایر بخش‌ها می‌گذارد.

  \item \textbf{توسعه تیمی و مقیاس‌پذیری:}
  در تیم‌های بزرگ، افراد/تیم‌های مختلف روی ویژگی‌های مختلف کار می‌کنند. تفکیک به ماژول‌ها احتمال Conflict و وابستگی‌های ناخواسته را کاهش می‌دهد.

  \item \textbf{\lr{Reusability}:}
  ماژول‌های عمومی مانند \lr{core:network}، \lr{core:database} یا \lr{core:designsystem} می‌توانند در پروژه‌های دیگر هم استفاده شوند بدون آنکه به منطق یا UI یک ویژگی خاص وابسته باشند.
\end{itemize}

\subsectionaddtolist{۲. معماری این پروژه به کدام الگو نزدیک‌تر است؟}

\begin{itemize}
  \item \textbf{\lr{MVVM}:}
  لایه رابط کاربری (View) با \lr{Jetpack Compose} ساخته می‌شود و وضعیت کاربر (State) و منطق نمایش در \lr{ViewModel} نگهداری می‌گردد.


  \item \textbf{\lr{Repository Pattern} و نزدیک به \lr{Clean Architecture}:} پروژه بیشتر به MVVM نزدیک است، اما می‌توان Clean هم در آن مشاهده کرد:
  \begin{itemize}
    \item \textbf{\lr{UI Layer}:} نمایش داده و مدیریت تعاملات.
    \item \textbf{\lr{Domain Layer}:} استفاده از \lr{UseCase}ها برای منطق تجاری.
    \item \textbf{\lr{Data Layer}:} مدیریت منابع داده (شبکه/دیتابیس) از طریق \lr{Repository}.
  \end{itemize}
\end{itemize}

\subsectionaddtolist{۳. چرا معماری نرم‌افزار به پروژه‌های وب محدود نمی‌شود؟}

معماری نرم‌افزار به وب محدود نیست، چون «معماری» به نوع پلتفرم وابسته نیست؛ بلکه روشی برای {مدیریت پیچیدگی}، {کاهش هزینه تغییرات} و {افزایش کیفیت} در هر سیستم نرم‌افزاری است. هرجا که یک سیستم شامل اجزای متعدد، وابستگی‌های گوناگون و نیاز به توسعه و نگه‌داری بلندمدت داشته باشد، معماری اهمیت پیدا می‌کند. دلایل اصلی عبارت‌اند از:

\begin{itemize}
  \item {همه سیستم‌های نرم‌افزاری با پیچیدگی و تغییرات مواجه‌اند:}
  چه وب باشد، چه موبایل، چه دسکتاپ، چه بازی، چه سامانه‌های نهفته یا \lr{IoT}، همگی با نیازمندی‌های جدید، تغییر فناوری، و رشد تدریجی سیستم روبه‌رو می‌شوند. معماری کمک می‌کند این تغییرات کنترل‌پذیر و کم‌ریسک‌تر شوند.

  \item {\lr{Non-Functional Requirements} عمومی‌اند:}
  مفاهیمی مثل {کارایی}، {مقیاس‌پذیری}، {امنیت}، {قابلیت اطمینان}، {نگه‌داری‌پذیری} و {تست‌پذیری} فقط مخصوص وب نیستند و در همه نرم‌افزارها مطرح می‌شوند. معماری ابزار تصمیم‌گیری برای رسیدن به این ویژگی‌هاست.

  \item {تفکیک مسئولیت‌ها و مدیریت وابستگی‌ها در همه‌جا لازم است:}
  اصل \lr{Separation of Concerns} و جلوگیری از وابستگی‌های درهم‌تنیده برای هر پروژه‌ای ضروری است؛ زیرا بدون آن، تغییر یک بخش باعث شکستن بخش‌های دیگر می‌شود و هزینه توسعه به‌مرور افزایش می‌یابد.

\end{itemize}

\noindent
بنابراین معماری نرم‌افزار یک مفهوم مستقل از پلتفرم است و در هر پروژه‌ای که قرار است پایدار، قابل توسعه، قابل تست و قابل نگه‌داری باشد، نقش اساسی دارد.

$\\$
